\section{Noisy expanding experiments and a uniform critical window}\label{sec:v4-expanding}

The realization overhead in Theorem~\ref{thm:v4-orbit} is not always
necessary. We now identify a class in which symbolic closure yields
matching errors, including the acquisition noise and the renewal tail.
Let $b\ge2$ and $d\ge1$ be integers, let $\mathbb T^d=(\R/\mathbb Z)^d$,
and put
\begin{equation}
 T_bx=bx\pmod1,\qquad
 \Phi(x)=\bigl(\cos(2\pi x_j),\sin(2\pi x_j)\bigr)_{j=1}^d.
                                                               \label{eq:v4-torus-readout}
\end{equation}
At a renewal $U$ is Haar uniform and the observation is
\begin{equation}
 Y=U+\sigma G\pmod1,\qquad G\sim N(0,I_d),\quad \sigma\ge0.
                                                               \label{eq:v4-wrapped-acquisition}
\end{equation}
The noise, the fresh states, and the renewal lengths are mutually
independent. The target is $\Phi(X_t)\in\R^{2d}$, not a discontinuous
choice of coordinates on the torus. Both average raw risk and average
excess risk have the meanings of Section~\ref{sec:v4-orbits}.

Haar invariance implies that $Y$ and $G$ are independent: conditional
on $G$, the translate of $U$ is Haar. Consequently the posterior law
of $U$ given $Y=y$ is that of $y-\sigma G$ modulo one, and
\begin{equation}
 m_k(y)=a_k\Phi(b^ky),\qquad
 a_k=\exp(-2\pi^2\sigma^2b^{2k}),\qquad
 B=d\sum_{k\ge0}w_k(1-a_k^2).                                  \label{eq:v4-explicit-posterior}
\end{equation}
This calculation uses the Gaussian characteristic function, so it
also specifies exactly which noise information is used by the decoder.
Write
\begin{equation}
 v_k=w_ka_k^2,\qquad
 \Psi_v(n)=\sum_{k<n}v_k b^{-2(n-k)}+\sum_{k\ge n}v_k,
 \quad n\ge0.                                                   \label{eq:v4-profile}
\end{equation}
For $v=w$, use the notation $\Psi_w$.
We impose the quantitative renewal condition
\begin{equation}
 W(n)\le A w_n\quad(n\ge0)                                     \label{eq:v4-residual-life}
\end{equation}
for some finite $A\ge1$. In terms of the cycle length this is a
uniform bound on the mean residual lifetime, conditional on survival.
It includes geometric renewals, mixtures with geometric parameters
bounded away from zero, and bounded renewal lengths. It does not
include arbitrary polynomial renewal tails.

Put $a=b^d$ and
\begin{equation}
 S_n=1+a+\cdots+a^n,\qquad
 n_M=\max\{n\ge0:S_n\le M\}.                                   \label{eq:v4-state-budget}
\end{equation}
The distinction between $a^n$ and $S_n$ counts, rather than suppresses,
the persistent information about remaining precision.

\begin{theorem}[Sharp posterior-orbit law]\label{thm:v4-noisy}
Under \eqref{eq:v4-residual-life}, for every $M\ge1$ and $\sigma\ge0$,
\begin{align}
 \frac{\Psi_v(n_M)}{8Ab^6}
 &\le E_M^{\rm av}\le \frac{d\pi^2}{3}\Psi_v(n_M),                \label{eq:v4-sharp-excess}\\
 B+\frac{\Psi_v(n_M)}{8Ab^6}
 &\le R_M^{\rm av}\le B+\frac{d\pi^2}{3}\Psi_v(n_M).             \label{eq:v4-sharp-raw}
\end{align}
The lower bounds hold for all randomized, nonstationary machines.
The upper bounds are attained by the finite suffix machine described
below. Its exact average raw risk with $S_n$ states is
\begin{equation}
 U_{n,\sigma}
 =B+d\sum_{k<n}v_k
       \left[1-\operatorname{sinc}^2(\pi b^{-(n-k)})\right]
       +d\sum_{k\ge n}v_k,                                     \label{eq:v4-exact-suffix}
\end{equation}
where $\operatorname{sinc}(z)=\sin(z)/z$ and $\operatorname{sinc}(0)=1$.
For geometric renewals, \eqref{eq:v4-sharp-raw} also holds for
$R_M^{\rm sup}:=\inf_{\mathcal A}\sup_{t\ge0}
\E\|\Phi(X_t)-\widehat f_t\|^2$, with an initial acquisition at zero.
The supremum raw risk of the displayed suffix machine is exactly
$U_{n,\sigma}$.
\end{theorem}

\begin{lemma}[Orbit separation and small balls]\label{lem:v4-smallballs}
Suppose $v_k=w_ka_k^2$ with $a_k$ nonincreasing and
\eqref{eq:v4-residual-life} holds. For Haar $Y$ the quantization error
of $(\sqrt{v_k}\Phi(b^kY))_{k\ge0}$ satisfies
\begin{equation}
 e_M^2\ge \frac{\Psi_v(n_M)}{8Ab^6}.                             \label{eq:v4-orbit-smallball}
\end{equation}
\end{lemma}
\begin{proof}
First, $\sum_{k\ge n}v_k\le A v_n$. Also $\Psi_v$ is nonincreasing,
and
\begin{equation}
 \Psi_v(n+1)\ge b^{-2}\Psi_v(n).                                \label{eq:v4-profile-step}
\end{equation}
These follow directly from \eqref{eq:v4-profile} and monotonicity of
$a_k$.
For $x,y\in\mathbb T^d$, write
$\delta=\max_j\operatorname{dist}_{\mathbb T}(x_j,y_j)\in[0,1/2]$.
If $\delta>0$, set $K=\lfloor\log_b(1/(2\delta))\rfloor$.
In a coordinate attaining $\delta$, for $0\le k\le K$ the unreduced
difference has modulus $b^k\delta\le1/2$. Hence
$|\sin(\pi b^k\delta)|\ge2b^k\delta$, and the squared orbit distance
$D^2(x,y)$ obeys
\begin{align}
 D^2(x,y)
 &=4\sum_{k\ge0}v_k\sum_{j=1}^d
          \sin^2(\pi b^k(x_j-y_j))\nonumber\\
 &\ge16\delta^2\sum_{k=0}^Kv_kb^{2k}
 \ge \frac4{b^2}\sum_{k=0}^Kv_kb^{-2(K-k)}
 \ge\frac4{Ab^2}\Psi_v(K).                                    \label{eq:v4-separated-orbits}
\end{align}
The middle inequality uses $\delta>b^{-(K+1)}/2$; the last uses the
residual-tail bound and $A\ge1$.

Let $h=[2(2M)^{1/d}]^{-1}$ and
$K_h=\lfloor\log_b(1/(2h))\rfloor=\lfloor\log_a(2M)\rfloor$.
Since $M<S_{n_M+1}$ and $a\ge2$, one has $K_h\le n_M+2$.
Whenever $\delta\ge h$, the associated $K\le K_h$; thus
\eqref{eq:v4-profile-step} and \eqref{eq:v4-separated-orbits} give
\[
 D^2(x,y)\ge\gamma\Psi_v(n_M),\qquad \gamma=\frac4{Ab^6}.
\]
If $\Psi_v(n_M)=0$, the conclusion is immediate. Otherwise take a
Hilbert ball of radius $r$, with
$r^2=\gamma\Psi_v(n_M)/16$. Two preimages in the same ball have
orbit distance at most $2r$, so their torus sup-distance is less than
$h$. Fixing one preimage, all preimages lie in a torus cube of side
$2h$, whose Haar measure is $1/(2M)$. This argument applies even when
the centre of the Hilbert ball is not on the orbit image.
The union of $M$ such balls has preimage measure at most $1/2$.
Every $M$-centre quantizer therefore has mean squared error at least
$r^2/2=\Psi_v(n_M)/(8Ab^6)$.
\end{proof}

\begin{proof}[Proof of Theorem~\ref{thm:v4-noisy}]
The lower bound follows from Theorem~\ref{thm:v4-orbit},
\eqref{eq:v4-explicit-posterior}, and
Lemma~\ref{lem:v4-smallballs}. The orthogonality identity then adds
$B$ to the excess risk.

For the upper bound use the alphabet $\{0,\ldots,b-1\}^d$.
The states are all words of lengths $0,\ldots,n$, exactly $S_n$ states.
At a renewal retain the first $n$ base-$b$ digits of each coordinate
of $Y$, grouped into one word. Use the terminating expansion at
ambiguous endpoints; those endpoints have Haar measure zero.
At a nonrenewal delete the first letter, leaving the empty word fixed.
For a nonempty word $w$ of length $l$, let $c(w)$ be the midpoint of
its $b$-adic cube, and decode it as
\[
 a_{n-l}\operatorname{sinc}(\pi b^{-l})\Phi(c(w)).
\]
Decode the empty word as zero. The age before exhaustion is $n-l$,
so there is no separate age register. The update does not use the
renewal law.

Conditional on this suffix at age $k<n$, the vector $b^kY$ is
uniform on its cube. Averaging each trigonometric coordinate over
an interval gives its midpoint value times
$\operatorname{sinc}(\pi b^{-(n-k)})$. The displayed output is
therefore the conditional mean of $\Phi(b^kU)$ under the suffix.
Its raw squared error is
\[
 d\left[1-a_k^2\operatorname{sinc}^2(\pi b^{-(n-k)})\right].
\]
After exhaustion the raw error is $d$. Renewal reward proves
\eqref{eq:v4-exact-suffix}. For $0\le t\le1/2$,
\begin{equation}
 \frac43t^2\le1-\operatorname{sinc}^2(\pi t)
                 \le\frac{\pi^2}{3}t^2.                        \label{eq:v4-sinc-bounds}
\end{equation}
Indeed, if $V,V'$ are independent uniform variables on an interval
of length $t$, the middle expression is
$2\E\sin^2(\pi(V-V'))$. Use
$2|u|\le|\sin(\pi u)|\le\pi|u|$ for $|u|\le1/2$ and
$\E(V-V')^2=t^2/6$. This proves the average upper bounds,
including $n=0$.

For geometric renewals put $q=1-p$. The raw age-loss sequence just
computed is nondecreasing: $a_k$ decreases and, as $k$ increases
before exhaustion, the sinc factor decreases as well. The age at
time $t$ has the law of $\min(G_0,t)$, where
$\Pp(G_0=k)=pq^k$. Therefore its expected raw loss increases to
$\sum_kpq^kd_k=U_{n,\sigma}$ and never exceeds this number.
The lower bound for supremum loss follows from the average converse.
\end{proof}

\begin{corollary}[Joint statistical and persistent-state law]\label{cor:v4-noise-memory}
Under the hypotheses of Theorem~\ref{thm:v4-noisy},
\begin{equation}
 R_M^{\rm av}\asymp_{b,d,A} B+\Psi_w(n_M).                       \label{eq:v4-noise-plus-memory}
\end{equation}
The same holds for raw supremum risk under geometric renewals.
Moreover a suffix decoder which does not know $\sigma$ or the
renewal law attains this order of raw risk.
\end{corollary}
\begin{proof}
Every coefficient in $\Psi_v$ is between zero and one. Thus
\[
 0\le\Psi_w(n)-\Psi_v(n)
 \le\sum_k w_k(1-a_k^2)=B/d.
\]
This and \eqref{eq:v4-sharp-raw} prove the equivalence.
For the noise-independent decoder use
$\operatorname{sinc}(\pi b^{-l})\Phi(c(w))$, omitting $a_{n-l}$.
Compare with the hypothetical target $\Phi(b^kY)$.
At age $k$,
\[
 \E\|\Phi(b^kU)-\Phi(b^kY)\|^2
 =2d(1-a_k)\le2d(1-a_k^2).
\]
The squared triangle inequality and the noiseless suffix bound give
average raw risk at most $4B+(2d\pi^2/3)\Psi_w(n)$.
For geometric renewals the same bound holds for supremum loss:
both the noise-discrepancy sequence $1-a_k^2$ and the noiseless
suffix raw-loss sequence are nondecreasing, so each time expectation
is bounded by its stationary-age average. The machine uses only
$b,d,n$ and the observed digits.
\end{proof}

\subsection{Geometric renewals and the critical window}
For $w_k=pq^k$, $p=1-q$, the noiseless profile is
\begin{equation}
 \Psi_q(n)
 =p b^{-2n}\sum_{k=0}^{n-1}(b^2q)^k+q^n.                        \label{eq:v4-geometric-profile}
\end{equation}
Theorem~\ref{thm:v4-noisy} is uniform for $q$ in any fixed compact
subset of $(0,1)$, since $A=1/p$ is then bounded. In particular it
remains valid when $q$ and $\sigma$ depend on $M$.

\begin{corollary}[Phase regimes and crossover]\label{cor:v4-critical}
For each fixed $0<q<1$, at $\sigma=0$ the optimal average and
supremum raw risks have orders
\begin{equation}
 \begin{cases}
 M^{-2/d},&q<b^{-2},\\
 (1+\log M)M^{-2/d},&q=b^{-2},\\
 M^{-\log(1/q)/(d\log b)},&q>b^{-2}.
 \end{cases}                                                    \label{eq:v4-three-regimes}
\end{equation}
Let $q_n=b^{-2}\exp(x/n)$ and
$H(x)=(e^x-1)/x$, with $H(0)=1$. Uniformly for $x$ in bounded
intervals,
\begin{equation}
 \frac{b^{2n}}n\Psi_{q_n}(n)
 \longrightarrow (1-b^{-2})H(x).                               \label{eq:v4-critical-window}
\end{equation}
For the explicit $S_n$-state suffix machine, its exact noiseless raw
risk satisfies the stronger coefficient limit
\begin{equation}
 \frac{b^{2n}}n U_{n,0}
 \longrightarrow \frac{d\pi^2}{3}(1-b^{-2})H(x).                 \label{eq:v4-suffix-crossover}
\end{equation}
The optimal risks are bounded above and below by constant multiples
of the profile in \eqref{eq:v4-critical-window}, uniformly in this
window; no equality of their leading constants is asserted.
\end{corollary}
\begin{proof}
Since $S_n\asymp_b b^{dn}$, evaluating the geometric sum in
\eqref{eq:v4-geometric-profile} proves
\eqref{eq:v4-three-regimes}. In the critical window,
\[
 \frac{b^{2n}}n\Psi_{q_n}(n)
 =(1-q_n)\frac1n\sum_{k=0}^{n-1}e^{xk/n}+\frac{e^x}{n}.
\]
The Riemann sum tends uniformly on bounded $x$-intervals to
$\int_0^1e^{xt}\,dt=H(x)$.
For the exact suffix risk use
$1-\operatorname{sinc}^2(\pi t)=(\pi^2/3)t^2+O(t^4)$,
uniformly for $0\le t\le1/2$. After multiplication by $b^{2n}$,
the error sum is bounded by a constant times
\[
 \sum_{k<n}e^{xk/n}b^{-2(n-k)},
\]
which is uniformly bounded on bounded $x$-intervals.
The exhaustion term $db^{2n}q_n^n=de^x$ is bounded as well.
Dividing by $n$ proves \eqref{eq:v4-suffix-crossover}.
The final assertion is exactly the uniform two-sided bound of
Theorem~\ref{thm:v4-noisy}.
\end{proof}

For completeness, the noise floor has the same transition. For
$0<\sigma\le1$, set
$n_\sigma=\lfloor\log_b(1/\sigma)\rfloor$; for $\sigma>1$ set
$n_\sigma=0$. From $1-e^{-4\pi^2z^2}\asymp\min(1,z^2)$ and
$b^{-(n_\sigma+1)}<\sigma\le b^{-n_\sigma}$ one obtains
\begin{equation}
 B\asymp_{b,d}\Psi_q(n_\sigma),\qquad
 R_M\asymp_{b,d,q} \Psi_q(n_\sigma)+\Psi_q(n_M).                 \label{eq:v4-joint-profile}
\end{equation}
At $\sigma=0$ the first term is interpreted as zero. The comparison
in the second formula is uniform on compact $q$-intervals. Thus the
small-noise orders of $B$ are respectively $\sigma^2$,
$\sigma^2(1+\log(1/\sigma))$, and
$\sigma^{\log(1/q)/\log b}$ in the three regimes. This describes
filtering uncertainty and expanding transport in the same experiment,
including joint limits, rather than adding an unrelated noisy model.

\subsection{Nonconstant slopes and other readouts}
The preceding law is not confined to a constant-slope coordinate
presentation. The next statement records both a genuine robustness
class and its precise boundary.

\begin{proposition}[Transport under conjugacy and metric readouts]\label{prop:v4-conjugacy}
Let $h:\mathbb T^d\to E'$ be a Borel isomorphism. Transport the reset
law, observation \eqref{eq:v4-wrapped-acquisition}, dynamics, and
readout by $h$; thus the new observation is $h(Y)$, the dynamics are
$hT_bh^{-1}$, and the readout is $\Phi h^{-1}$. All the preceding
finite-state risks are unchanged.
Suppose instead that $E'$ is compact and the new continuous readout
$g$ satisfies, for some $0<c\le C<\infty$,
\begin{equation}
 c\|\Phi(x)-\Phi(y)\|
 \le\|g(hx)-g(hy)\|
 \le C\|\Phi(x)-\Phi(y)\|.                                    \label{eq:v4-readout-comparison}
\end{equation}
Then its optimal \emph{raw} squared risks, for the same observations
and state budgets, lie between $c^2/4$ and $4C^2$ times those for
$\Phi$. This assertion concerns raw risk, not equality of posterior
excess risk for nonlinear readouts.
\end{proposition}
\begin{proof}
Apply $h^{-1}$ to each acquired observation within the transient
update computation. This and its inverse transform machines without
changing a persistent state. For the readout assertion, replace any
output $z$ for $\Phi$ by a nearest point $\Phi(x_z)$ on its compact
image, choosing measurably. For a true state $x$,
$\|\Phi(x_z)-\Phi(x)\|\le2\|z-\Phi(x)\|$.
Output $g(hx_z)$ and use the upper inequality in
\eqref{eq:v4-readout-comparison}. The inverse construction projects
onto the compact image of $g$ and uses the lower inequality.
Both transformations act only on the decoder, use no extra states,
and are valid conditionally on any coding randomness.
\end{proof}

For example, in one dimension take
\[
 h_\varepsilon(x)=x+\frac{\varepsilon}{2\pi}\sin(2\pi x)\pmod1,
 \qquad |\varepsilon|<\frac{b-1}{b+1}.
\]
This is a smooth circle diffeomorphism and the conjugate map has
nonconstant derivative
\[
 (h_\varepsilon T_bh_\varepsilon^{-1})'(h_\varepsilon x)
 =b\frac{1+\varepsilon\cos(2\pi bx)}{1+\varepsilon\cos(2\pi x)}>1.
\]
Coordinatewise products give higher-dimensional examples, with
nonuniform physical reset density and transported nonadditive
observation noise. The ordinary circle embedding of $h_\varepsilon x$
satisfies \eqref{eq:v4-readout-comparison}. These are nonconstant-slope
expanding experiments, but their Lyapunov exponent remains $\log b$
by conjugacy. We do not identify their law with a pressure theorem
for arbitrary nonconjugate expanding maps. The general posterior-orbit
theorem still applies to such maps; its quantization invariant would
have to be evaluated separately.
